%% This document is created by 
%%  Dr. Putu Harry Gunawan
%% and edited by 
%%  Gia S. Wulandari (2023)
%% Template untuk Proposal TA 1 dan TA
%% Template ini digunakan untuk penulisan proposal TA 1 atau TA Fakultas Ilmu Komputer & Informatika, Universitas Terakreditasi.
%%  New Update 28 Pebruari 2017 : Bookmarks pdf, Pembimbing hanya satu

\documentclass[a4paper,12pt,oneside]{book}
\usepackage[utf8]{inputenc}
\usepackage{sectsty}
\usepackage{graphicx}
\usepackage{epstopdf}
\usepackage{algorithm}
\usepackage{algpseudocode}
\usepackage{array}
\usepackage[table]{xcolor}
\usepackage{anysize}
\usepackage{amsmath}
\usepackage{amssymb}
\usepackage[bahasa]{babel}
\usepackage{indentfirst} %Spasi untuk paragraf pertama
\usepackage{geometry}
\usepackage{multirow}% http://ctan.org/pkg/multirow
\usepackage{hhline}% http://ctan.org/pkg/hhline
\marginsize{4cm}{3cm}{3cm}{3cm} %{left}{right}{top}{bottom}
\usepackage[compact]{titlesec} 
\usepackage{etoolbox}
\usepackage{natbib} %for Harvard syle
\usepackage[bookmarks,hypertexnames=false,debug]{hyperref}
%\usepackage[pageanchor]{hyperref}
\usepackage{bookmark}


\makeatletter
\patchcmd{\ttlh@hang}{\parindent\z@}{\parindent\z@\leavevmode}{}{}
\patchcmd{\ttlh@hang}{\noindent}{}{}{}
\makeatother

\chapterfont{\centering}
\newcommand{\bigsize}{\fontsize{16pt}{14pt}\selectfont}
\chapterfont{\centering\bigsize\bfseries}
\sectionfont{\large\bfseries}
\usepackage{tikz}
\usetikzlibrary{shapes.geometric, arrows}
%\renewcommand{\chaptertitle}{BAB}
\renewcommand{\thechapter}{\Roman{chapter}}
\renewcommand\thesection{\arabic{chapter}.\arabic{section}}
\renewcommand\thesubsection{\thesection.\arabic{subsection}}
\renewcommand{\theequation}{\arabic{chapter}.\arabic{equation}}
\renewcommand{\thefigure}{\arabic{chapter}.\arabic{figure}}
\renewcommand{\thetable}{\arabic{chapter}.\arabic{table}}

\renewcommand\bibname{Daftar Pustaka}
\addto{\captionsbahasa}{\renewcommand{\bibname}{Daftar Pustaka}}
\usepackage{fancyhdr}
\pagestyle{fancy}
\lhead{}
\chead{}
\rhead{}
\lfoot{}
\cfoot{\thepage}
\rfoot{}
\renewcommand{\headrulewidth}{0pt}

\makeatletter

%%%%%%%%%%%%%%%%%%%%%%%%%%%%%%%%%%%%%%%%%%%%%%%%%%%%%%%%%%%%
%
%  Berikut adalah data-data yang wajib diisi oleh mahasiswa
%
%%%%%%%%%%%%%%%%%%%%%%%%%%%%%%%%%%%%%%%%%%%%%%%%%%%%%%%%%%%%

\title{Judul Proposal TA Singkat dan Spesifik, Tetapi Cukup Jelas Memberi Gambaran Mengenai TA yang Diusulkan}\let\Title\@title   %Judul dalam bahasa Indonesia

\newcommand{\EngTitle}{Judul Proposal TA Dalam Bahasa Inggris}  %Judul dalam bahasa Inggris

\author{Mulyani}  \let\Author\@author  %Nama mhs
\newcommand{\NIM}{(NIM)} %Jangan lupa hapus tanda ( dan )
\newcommand{\Prodi}{Informatika}
\newcommand{\KK}{(Kelompok Keahlian)} %UNTUK TA, Jangan lupa hapus tanda ( dan )
\newcommand{\Gelar}{S.Kom.} % UNTUK TA
\date{\the\year}           \let\Date\@date %Masukkan hanya tahun saja
\newcommand{\Tanggal}{\the\day} % Tanggal Pengesahan
\newcommand{\Bulan}{Oktober} % Bulan Pengesahan
\newcommand{\PembimbingSatu}{(Dr. Calon Pembimbing 1)}
\newcommand{\NIPPembimbingSatu}{123456}
\newcommand{\PembimbingDua}{(Dr. Calon Pembimbing 2)}
\newcommand{\NIPPembimbingDua}{123456}
\newcommand{\Kaprodi}{Nama Pembimbing}
\newcommand{\NIPKaprodi}{00760045}
\newif\ifPembimbingHanyaSatu

%%%% WARNING kode berikut ini diaktifkan Jika pembimbing hanya satu %%%%

%\PembimbingHanyaSatutrue

%%%%%%%%%%%%%%%%%%%%%%%%%%%%%%%%%%%%%%%%%%%%%%%%%%%%%%%%%%%%%%%%%%

\makeatother
\linespread{1}


\begin{document}
%%%%%%%%%%%%%%%%%%%%%%%%%%%%%%%%%%%%%%%%%%%%%%%%%%%%%%%%%%
% Dimulai dengan cover, lembar persetujuan, Abstrak, dll
%%%%%%%%%%%%%%%%%%%%%%%%%%%%%%%%%%%%%%%%%%%%%%%%%%%%%%%%%%%
\pagenumbering{roman} 
\begin{titlepage}
\thispagestyle{empty}
%\vspace*{0.7cm}
\pdfbookmark{Cover}{ }
\include{Cover}
\pagebreak
\thispagestyle{empty}
\pdfbookmark{Lembar Persetujuan}{ }
\include{Lembar-Persetujuan}
\pagebreak
\end{titlepage}
\phantomsection
\addcontentsline{toc}{chapter}{Abstrak}
%\pdfbookmark{\abstractname}{Abstrak Indonesia}
\include{Abstrak-Indo}
\cleardoublepage
\phantomsection
\addcontentsline{toc}{chapter}{Daftar Isi}
%\pdfbookmark{\contentsname}{Contents}
\tableofcontents

%%%%%%%%%%%%%%%%%%%%%%%%%%%%%%%%%%%%%%%%%%%%%%%%%
% Bab bab penulisan
%%%%%%%%%%%%%%%%%%%%%%%%%%%%%%%%%%%%%%%%%%%%%%%%
\cleardoublepage
\pagenumbering{arabic}
\chapter{Pendahuluan}
Bagian pendahuluan memuat beberapa substansi sebagai berikut:

\section{Latar Belakang}
TA dilakukan untuk menjawab keingintahuan mahasiswa mengenai suatu gejala/konsep/dugaan. Kemukakan argumentasi pentingnya dilakukan pengerjaan TA yang diusulkan tersebut dengan menyampaikan hasil-hasil penelitian sebelumnya yang ada pada referensi. Secara umum latar belakang berisi:
• Alasan kenapa kasus/masalah/fenomena tersebut diambil sebagai bahan kajian.
• Apakah ada sebuah konsep baru sebagai hasil penelitian?
• Gap antara kondisi saat ini dengan kondisi yang akan datang (diharapkan).
Uraikan proses-proses yang dilakukan dalam mengidentifikasi masalah yang akan dicari solusinya. Latar belakang minimal 1 halaman.

\section{Perumusan Masalah}
Uraikan permasalahan yang akan dibahas dalam TA dengan mengacu pada latar belakang yang telah disampaikan dan hasil penelitian terdahulu (bila ada). Dalam perumusan masalah dapat dijelaskan definisi, asumsi, dan lingkup yang menjadi batasan TA. Uraian perumusan masalah tidak harus dalam bentuk pertanyaan

\section{Tujuan}
Berikan pernyataan singkat mengenai tujuan TA. Tujuan dapat berupa menguraikan, menerangkan, membuktikan atau menerapkan suatu gejala/konsep/dugaan, atau membuat suatu model. Rumuskan tujuan yang akan dicapai secara spesifik yang merupakan kondisi baru yang diharapkan terwujud setelah TA selesai. Tujuan harus jelas dan dapat diukur, serta berhubungan dengan perumusan masalah yang telah dituliskan. Batasan masalah pun dapat ditambahkan pada bagian ini.
        
\section{Hipotesis (opsional)}
Merupakan jawaban sementara atas permasalahan yang akan dibahas. Hipotesis memuat penjelasan mengenai metode/konsep yang akan digunakan untuk memecahkan masalah serta alasan pemilihan metode/konsep tersebut. Selain itu, dalam hipotesis dimunculkan perbedaan antar metode/konsep yang digunakan dengan metode/konsep terdahulu. Hipotesis juga berisikan perkiraan hasil dari rencana solusi yang akan dilakukan.


\section{Rencana Kegiatan}
Rencana kegiatan adalah penjelasan mengenai rencana langkah-langkah yang akan dilakukan dalam pengerjaan Tugas Akhir yang memuat: kajian pustaka, cara pengumpulan data (kualitatif, kuantitatif), rancangan penelitian (mencakup prosedur penelitian dan perancangan sistem), cara menguji hasil penelitian (cara penafsiran dan penyimpulan hasil penelitian).


\section{Jadwal Kegiatan}

Jadwal pelaksanaan dibuat berdasarkan metodologi penyelesaian masalah yang digunakan. Bar-chart bisa dibuat per bulan atau per minggu. Dibawah ini adalah merupakan contoh bar-chart:

\begin{table}[h!]
  \centering
    \caption{Jadwal kegiatan proposal tugas akhir}
  \label{Novella}
  \begin{tabular}{|c|m{2.5cm}|m{0.01cm}|m{0.01cm}|m{0.01cm}|m{0.01cm}|m{0.01cm}|m{0.01cm}|m{0.01cm}|m{0.01cm}|m{0.01cm}|m{0.01cm}|m{0.01cm}|m{0.01cm}|m{0.01cm}|m{0.01cm}|m{0.01cm}|m{0.01cm}|m{0.01cm}|m{0.01cm}|m{0.01cm}|m{0.01cm}|m{0.01cm}|m{0.01cm}|m{0.01cm}|m{0.01cm}|}
    \hline
    \multirow{2}{*}{\textbf{No}} & \multirow{2}{*}{\textbf{Kegiatan}} & \multicolumn{24}{|c|}{\textbf{Bulan ke-}} \\
    \hhline{~~------------------------}
    {} & {} & \multicolumn{4}{|c|}{\textbf{1}} & \multicolumn{4}{|c|}{\textbf{2}} & \multicolumn{4}{|c|}{\textbf{3}} & \multicolumn{4}{|c|}{\textbf{4}} & \multicolumn{4}{|c|}{\textbf{5}} & \multicolumn{4}{|c|}{\textbf{6}}\\
    \hline
    1 & Studi Literatur & \cellcolor{blue!25} & \cellcolor{blue!25} & \cellcolor{blue!25} & \cellcolor{blue!25}& \cellcolor{blue!25} & \cellcolor{blue!25} & \cellcolor{blue!25} & \cellcolor{blue!25}& \cellcolor{blue!25} & \cellcolor{blue!25} & \cellcolor{blue!25} & \cellcolor{blue!25}& \cellcolor{blue!25} & \cellcolor{blue!25} & \cellcolor{blue!25} & \cellcolor{blue!25}& \cellcolor{blue!25} & \cellcolor{blue!25} & \cellcolor{blue!25} & \cellcolor{blue!25}& \cellcolor{blue!25} & \cellcolor{blue!25} & \cellcolor{blue!25} & \cellcolor{blue!25}\\
    \hline
    2 & Pengumpulan Data & \cellcolor{blue!25} & \cellcolor{blue!25} & \cellcolor{blue!25} & \cellcolor{blue!25} & {} & {} & {} & {} & {} & {} & {} & {}& {} & {} & {} & {}& {} & {} & {} & {}& {} & {} & {} & {}\\
    \hline
    3 & Analisis dan Perancangan Sistem &  {} & {} & {} & {}  & \cellcolor{blue!25} & \cellcolor{blue!25} & \cellcolor{blue!25} & \cellcolor{blue!25} & \cellcolor{blue!25} & \cellcolor{blue!25} & \cellcolor{blue!25} & \cellcolor{blue!25} & {} & {} & {} & {}& {} & {} & {} & {}& {} & {} & {} & {}\\
    \hline
    4 & Implementasi Sistem &  {} & {} & {} & {} & {} & {} & {} & {}& \cellcolor{blue!25} & \cellcolor{blue!25} & \cellcolor{blue!25} & \cellcolor{blue!25} & \cellcolor{blue!25} & \cellcolor{blue!25} & \cellcolor{blue!25} & \cellcolor{blue!25} & {} & {} & {} & {}& {} & {} & {} & {}\\
    \hline
    5 & Analisa Hasil Implementasi &  {} & {} & {} & {} & {} & {} & {} & {}& {} & {} & {} & {} & \cellcolor{blue!25} & \cellcolor{blue!25} & \cellcolor{blue!25} & \cellcolor{blue!25} & \cellcolor{blue!25} & \cellcolor{blue!25} & \cellcolor{blue!25} & \cellcolor{blue!25} & {} & {} & {} & {}\\
    \hline
    6 & Penulisan Laporan & {} & {} & {} & {} & \cellcolor{blue!25} & \cellcolor{blue!25} & \cellcolor{blue!25} & \cellcolor{blue!25}& \cellcolor{blue!25} & \cellcolor{blue!25} & \cellcolor{blue!25} & \cellcolor{blue!25}& \cellcolor{blue!25} & \cellcolor{blue!25} & \cellcolor{blue!25} & \cellcolor{blue!25}& \cellcolor{blue!25} & \cellcolor{blue!25} & \cellcolor{blue!25} & \cellcolor{blue!25}& \cellcolor{blue!25} & \cellcolor{blue!25} & \cellcolor{blue!25} & \cellcolor{blue!25}\\
    \hline
  \end{tabular}

\end{table}
\newpage


%
\include{Kajian-Pustaka}
%
\include{Metodologi}
%

%%%%%%%%%%%%%%%%%%%%%%%%%%%%%%%%%%%%%%%%%%%%%%%%%%%%%%%%%%%%%%%
% Berikut adalah untuk membuat Daftar pustaka
% Style daftar pustaka dapat disesuaikan dengan mengubah 
%\bibliographystyle{kode}
% dengan kode = acm maka hasinya contoh [1]
% dengan kode = agsm maka hasilnya Harvard style (Sukimin, 2017)
%%%%%%%%%%%%%%%%%%%%%%%%%%%%%%%%%%%%%%%%%%%%%%%%%%%%%%%%%%%%%%%%
\cleardoublepage
\addcontentsline{toc}{chapter}{Daftar Pustaka}

\bibliographystyle{plain} %harvard style
\bibliography{References}
%
%\pagebreak
%%%%%%%%%%%%%%%%%%%%%%%%%%%%%%%%%%%%%%%%%%%%%%%%%%%%%%%%%%%%%%%
% Berikut untuk lampiran
%%%%%%%%%%%%%%%%%%%%%%%%%%%%%%%%%%%%%%%%%%%%%%%%%%%%%%%%%%%%%%
\cleardoublepage
\addcontentsline{toc}{chapter}{Lampiran}
\include{Lampiran}
\end{document}
